\documentclass[12pt,a4paper]{article}

\usepackage[T2A]{fontenc}
\usepackage[utf8]{inputenc}
\usepackage[russian]{babel}
\usepackage{geometry}
\usepackage{array}
\usepackage{longtable}
\usepackage{tabularx}
\usepackage{booktabs}
\usepackage{hyperref}
\usepackage{enumitem}
\usepackage{setspace}
\usepackage{titlesec}

\geometry{margin=2.2cm}
\setstretch{1.15}
\hypersetup{
    colorlinks=true,
    linkcolor=black,
    urlcolor=blue
}
\setlist[itemize]{leftmargin=1.8em}
\setlist[enumerate]{leftmargin=2em}
\setlength{\parindent}{1.25cm}
\setlength{\parskip}{0.35em}

\newcommand{\ReportTitle}{Отчёт по практике: анализ истории разработки llm-tools-platform}
\newcommand{\ReportSubtitle}{Дневник работ и аналитический обзор по полной Git-истории проекта}
\newcommand{\ReportAuthor}{Заполнить ФИО}
\newcommand{\ReportGroup}{Заполнить группу}
\newcommand{\ReportSupervisor}{Заполнить руководителя}
\newcommand{\ReportUniversity}{Заполнить вуз}
\newcommand{\ReportDepartment}{Заполнить кафедру}
\newcommand{\ReportCity}{Заполнить город}
\newcommand{\ReportYear}{2026}

\titleformat{\section}{\large\bfseries}{\thesection.}{0.6em}{}
\titleformat{\subsection}{\normalsize\bfseries}{\thesubsection.}{0.6em}{}

\begin{document}

\begin{titlepage}
\begin{center}
{\small \ReportUniversity \par}
{\small \ReportDepartment \par}
\vfill
{\Large \textbf{\ReportTitle} \par}
\vspace{0.8cm}
{\large \ReportSubtitle \par}
\vfill
\begin{flushright}
\begin{tabular}{rl}
Студент: & \ReportAuthor \\
Группа: & \ReportGroup \\
Руководитель: & \ReportSupervisor \\
\end{tabular}
\end{flushright}
\vfill
{\ReportCity\ --- \ReportYear \par}
\end{center}
\end{titlepage}

\tableofcontents
\newpage

\section{Введение}

В доступной истории проекта \texttt{llm-tools-platform} зафиксировано \textbf{256 коммитов} в диапазоне от \texttt{2025-01-01} до \texttt{2026-04-12}. При этом реальная рабочая линия начинается с \texttt{2026-01-05}: более ранний коммит выглядит как стартовый шаблон, а основная инженерная активность появляется уже в январе 2026 года. Самым насыщенным месяцем оказался март 2026 года, когда в репозиторий вошло \textbf{150 коммитов}; именно в этот период проект сделал переход от набора отдельных функций к более зрелой backend-платформе с выраженным операционным контуром.

Цель отчёта --- не просто перечислить коммиты по датам, а показать, как по истории репозитория читается развитие системы: от ранней сборки каркаса через переход к \texttt{v3.0} и усиление \texttt{UMS}/\texttt{RAG} до текущей консолидации вокруг \texttt{Open WebUI-first}. Поэтому документ построен как гибрид дневника практики и аналитической записки: сначала фиксируется фактическая последовательность работ, затем объясняется, зачем эти шаги были нужны и к какому состоянию они привели.

Практический смысл такого формата в том, что история Git позволяет увидеть не только результат, но и сам ход инженерной мысли. По этой истории можно отделить быстрые экспериментальные шаги от осознанных архитектурных решений, понять, когда проект находился в режиме стабилизации, а когда --- в режиме ускоренного наращивания возможностей, а также выделить точки, где команда явно меняла приоритеты и целевую архитектуру.

\section{Методика анализа}

Для подготовки отчёта использовались только данные, которые доступны в самом репозитории на текущем состоянии ветки \texttt{refactor/openwebui-responsibility-split}. В выборку вошли:

\begin{itemize}
\item полная история всех локально доступных веток через \texttt{git log --all};
\item статистика авторов через \texttt{git shortlog -sne --all};
\item сравнение веток через \texttt{git rev-list --left-right --count};
\item фиксация тематических этапов по сообщениям коммитов и по зонам репозитория, которые менялись чаще всего.
\end{itemize}

Из анализа исключены незакоммиченные изменения рабочего дерева, потому что они не являются частью подтверждённой истории проекта. Отдельно был нормализован вклад основного автора, так как в истории используются несколько \texttt{git}-идентичностей с одними и теми же рабочими паттернами. Боты и вспомогательные агентные учётные записи рассматривались как отдельные участники, но в аналитической части трактуются как вспомогательный инструмент внутри основной линии разработки.

Для того чтобы отчёт не превращался в механическую распечатку \texttt{git log}, история была разложена на четыре содержательные фазы:

\begin{enumerate}
\item сборка основы и ранняя интеграция;
\item переход к \texttt{v3.0} и аппаратно-адаптивной архитектуре;
\item backend-first усиление и операционализация;
\item консолидация вокруг \texttt{Open WebUI-first}.
\end{enumerate}

Такой подход удобен сразу по двум причинам. Во-первых, он сохраняет дневниковую достоверность: по каждой активной дате видно, что именно делалось. Во-вторых, он даёт аналитическую рамку, в которой отдельные коммиты перестают быть случайным списком и складываются в понятные инженерные переходы.

\section{Общая картина истории проекта}

\begin{itemize}
\item Всего коммитов в анализируемой истории: 256.
\item Диапазон истории: 2025-01-01 -- 2026-04-12.
\item Рабочая линия проекта начинается с 2026-01-05, потому что более ранний коммит выглядит как шаблонный стартовый импорт.
\item Текущая рабочая ветка на момент подготовки отчёта: \texttt{refactor/openwebui-responsibility-split}.
\item Всего активных дат с коммитами: 38.
\end{itemize}

\begin{table}[h]
\centering
\begin{tabular}{lr}
\textbf{Автор} & \textbf{Коммитов} \\
\hline
senorfrancesco & 215 \\
Manus AI & 19 \\
copilot-swe-agent[bot] & 10 \\
gpt-engineer-app[bot] & 6 \\
anthropic-code-agent[bot] & 5 \\
\end{tabular}
\caption{Наиболее активные авторы в истории проекта}
\end{table}

\begin{table}[h]
\centering
\begin{tabular}{lr}
\textbf{Дата} & \textbf{Коммитов} \\
\hline
2026-01-06 & 36 \\
2026-03-18 & 28 \\
2026-03-13 & 26 \\
2026-03-19 & 15 \\
2026-03-04 & 12 \\
\end{tabular}
\caption{Самые насыщенные дни разработки}
\end{table}


Статистика хорошо показывает, что проект развивался неравномерно. Январь 2026 года --- это фаза быстрой сборки и экстренной стабилизации: много разнотипных коммитов, часть сообщений ещё нерегулярна, заметны развороты, повторные исправления и даже откаты. Февраль выглядит короче, но именно здесь начинается переход к более осмысленной версии \texttt{v3.0}. Март, напротив, показывает инженерную зрелость: в истории резко растёт доля \texttt{feat}, \texttt{docs}, \texttt{fix} и \texttt{test}, то есть новые возможности появляются вместе с документацией, проверками и инфраструктурой эксплуатации.

С точки зрения состава работ проект уже давно не сводится только к интерфейсу чата или разовой обработке документов. По коммитам видно развитие нескольких самостоятельных слоёв: orchestration-логики, работы моделей и \texttt{UMS}, поиска и знаний, операторского интерфейса, сценариев развертывания, а затем --- миграционного и \texttt{Open WebUI}-ориентированного контура. Именно поэтому для итогового анализа важен не только объём работ, но и их последовательность.

\section{Дневниковая часть: ход работ по активным датам}

В этом разделе сохранён формат дневника практики. Пустые календарные дни в таблицу не включались: оставлены только те даты, в которые история Git действительно фиксирует работу. Такой подход делает дневник компактнее, но при этом не теряет фактическую хронологию.

\input{generated/practice_history_timeline.tex}

\section{Аналитический обзор этапов}

Ниже собраны четыре укрупнённые фазы, внутри которых особенно хорошо видно, как менялись задачи проекта и инженерные приоритеты.

\subsection*{Фаза 1. Сборка основы и ранняя интеграция}
\textbf{Период:} 2026-01-05 -- 2026-01-09 \\
\textbf{Коммитов в фазе:} 53

Проект собирается из нескольких заготовок в единую систему: складываются backend, UI, потоковый режим, \texttt{UMS} и первые пользовательские сценарии. По коммитам доминируют общая стабилизация и организационные правки и модели, \texttt{UMS}, распределение ресурсов и наблюдаемость. Это был период шлифовки, когда устойчивость становилась важнее простого прироста функций.

\begin{itemize}
\item \texttt{2b2010c} (2026-01-05) --- merge with all proj
\item \texttt{d0c788d} (2026-01-07) --- feat: Integrate Open WebUI, fix UMS inference, add MD reporting
\item \texttt{e7e7645} (2026-01-06) --- Chore: Clean up warnings and modernize FastAPI events in UMS
\item \texttt{f79bddb} (2026-01-06) --- Feature: Add .env support and clean up GPU monitoring imports
\item \texttt{62ff5e8} (2026-01-06) --- Feature: End-to-end real-time streaming
\end{itemize}

\subsection*{Фаза 2. Переход к v3.0 и аппаратно-адаптивной архитектуре}
\textbf{Период:} 2026-02-24 -- 2026-03-05 \\
\textbf{Коммитов в фазе:} 33

Появляется более зрелый контур \texttt{v3.0}: \texttt{Chainlit}, \texttt{AdaptiveRAGPipeline}, аппаратное профилирование, новые workflow и тестовое покрытие. По коммитам доминируют поддержка и эволюция \texttt{Chainlit}-контура и общая стабилизация и организационные правки. Это был период ускоренного наращивания продукта и архитектуры.

\begin{itemize}
\item \texttt{ceae15c} (2026-02-26) --- feat: v3.0 Hardware-Adaptive Agentic System — profiler, tiered RAG, Chainlit UI
\item \texttt{adf9bbf} (2026-03-04) --- chore(repo): ignore local test logs and ad-hoc node bootstrap
\item \texttt{512fd97} (2026-03-04) --- chore(repo): remove obsolete verification script
\item \texttt{0f59b70} (2026-02-26) --- chore: update .gitignore and requirements.txt for v3.0
\item \texttt{5c62f9a} (2026-03-04) --- docs(repo): consolidate task docs and ignore local artifacts
\end{itemize}

\subsection*{Фаза 3. Backend-first усиление и операционализация}
\textbf{Период:} 2026-03-12 -- 2026-03-31 \\
\textbf{Коммитов в фазе:} 124

Команда переводит проект от набора функций к платформе: усиливаются \texttt{orchestrator}, \texttt{UMS}, реестр моделей, \texttt{operator UI}, офлайн-поставка и наблюдаемость. По коммитам доминируют рантайм, запуск, офлайн-развёртывание и инсталляция и документация и фиксация правил эксплуатации. Это был период ускоренного наращивания продукта и архитектуры.

\begin{itemize}
\item \texttt{85bd9cb} (2026-03-18) --- feat(model-registry): add canonical yaml registry
\item \texttt{fac57e5} (2026-03-29) --- feat(operator): add python control plane for operator ui
\item \texttt{5ac5f34} (2026-03-31) --- Fix compare workflow LLM queue discrepancy and JSON parsing
\item \texttt{8ef99cd} (2026-03-17) --- chore(config): add local hardware override defaults
\item \texttt{d91b8f1} (2026-03-16) --- chore(git): ignore local runtime artifacts
\end{itemize}

\subsection*{Фаза 4. Консолидация вокруг Open WebUI-first}
\textbf{Период:} 2026-04-01 -- 2026-04-12 \\
\textbf{Коммитов в фазе:} 44

Основной вектор смещается к \texttt{Open WebUI}: уточняется разделение ответственности между инструментами, bootstrap-процессом, рантаймом и миграционной документацией. По коммитам доминируют документация и фиксация правил эксплуатации и миграция и стабилизация \texttt{Open WebUI}. Это был период ускоренного наращивания продукта и архитектуры.

\begin{itemize}
\item \texttt{71c9c51} (2026-04-11) --- feat(openwebui): reconcile bootstrap drift and native function-calling defaults
\item \texttt{4a800b1} (2026-04-01) --- Fix model download scripts: Qwen3-14B from MaziyarPanahi, LaBSE from cointegrated
\item \texttt{c8e8671} (2026-04-01) --- chore: update package version lower bounds and build script version pins
\item \texttt{062d6e9} (2026-04-07) --- docs(repo): relocate forms and research artifacts
\item \texttt{02d83c8} (2026-04-01) --- feat(compare): implement B3.51g and B3.51h - summary-first bounded latency with fast/deep modes
\end{itemize}


\subsection{Что означают эти фазы на практике}

Первая фаза выглядит как период, когда команда быстро соединяла несколько ранее независимых заготовок в один рабочий контур. Именно поэтому в январской истории много разнородных исправлений: шла не косметическая доработка, а фактическое приведение системы к состоянию, в котором она вообще может устойчиво запускаться и демонстрировать целостный сценарий.

Вторая фаза уже заметно осмысленнее. Переход к \texttt{v3.0}, появление \texttt{Hardware-Adaptive Agentic System}, \texttt{AdaptiveRAGPipeline} и новых workflow означает, что проект начал рассматривать себя не как простой интерфейс к модели, а как платформу, которая должна учитывать тип задачи, состояние оборудования и разный характер документов.

Третья фаза --- ключевая с точки зрения инженерной зрелости. Здесь история начинает говорить языком платформенной разработки: появляются реестр моделей, контуры отказоустойчивости, политика распределения ресурсов, операторский интерфейс, офлайн-поставка, метрики, наблюдаемость. Это уже не просто накопление функций, а перевод системы в состояние, где её можно сопровождать, проверять и разворачивать в более строгом режиме.

Четвёртая фаза показывает, что команда перестала держать несколько равноправных пользовательских путей и начала явно консервировать целевую архитектуру вокруг \texttt{Open WebUI}. Это важный признак зрелости: на позднем этапе ценность даёт не только добавление функций, но и сокращение архитектурной двусмысленности.

\section{Анализ веток и линий разработки}

\begin{table}[h]
\centering
\begin{tabular}{p{0.25\textwidth} p{0.5\textwidth} r}
\textbf{Ветка} & \textbf{Смысл} & \textbf{vs dev} \\
\hline
\texttt{dev} & Основная интеграционная ветка, в которой сходятся рабочие срезы перед закреплением следующего состояния проекта. & +0 / -0 \\
\texttt{refactor/openwebui-responsibility-split} & Текущая активная линия, где проект закрепляет \texttt{Open WebUI} как основной интерфейс и разводит зоны ответственности между bootstrap, tool-контуром и рантаймом. & +14 / -0 \\
\texttt{release/v0.1} & Исторический релизный срез на этапе начала миграционного контура; полезен как точка сравнения перед дальнейшей консолидацией. & +0 / -19 \\
\texttt{v0.1\_history} & Дублирующий исторический срез той же точки, сохранённый как архивная ветка состояния. & +0 / -19 \\
\end{tabular}
\caption{Ключевые ветки, которые использованы в отчёте}
\end{table}

\paragraph{Точечные удалённые ветки.} Помимо основных линий, в истории есть короткие вспомогательные ветки под узкие задачи. Они важны как маркеры направления, но не образуют самостоятельную продуктовую линию.
\begin{itemize}
\item \texttt{origin/claude/compare-workflow-llm-discrepancy} --- Точечная ветка для корректировки compare-workflow и снижения расхождений в LLM-обработке.
\item \texttt{origin/claude/redesign-compare-bounded-latency} --- Точечная ветка для корректировки compare-workflow и снижения расхождений в LLM-обработке.
\item \texttt{origin/codex/orchestration-control-plane-snapshot} --- Снимок состояния оркестрационного контура в отдельный момент разработки.
\item \texttt{origin/copilot/check-pull-requests-for-v3-0} --- Служебная ветка для узкой правки зависимостей, очистки или проверки конкретного сценария.
\item \texttt{origin/copilot/fix-model-download-scripts} --- Служебная ветка для узкой правки зависимостей, очистки или проверки конкретного сценария.
\item \texttt{origin/copilot/remove-legacy-code-and-update-docs} --- Служебная ветка для узкой правки зависимостей, очистки или проверки конкретного сценария.
\item \texttt{origin/copilot/update-package-scripts-and-dependencies} --- Служебная ветка для узкой правки зависимостей, очистки или проверки конкретного сценария.
\end{itemize}


Если смотреть на ветки не формально, а по их смыслу, то картина получается достаточно стройной. \texttt{dev} выступает как рабочая интеграционная база. \texttt{release/v0.1} и \texttt{v0.1\_history} полезны не тем, что дают параллельную историю, а тем, что фиксируют состояние проекта на важном переходном моменте. Текущая ветка \texttt{refactor/openwebui-responsibility-split} показывает уже следующий шаг --- не просто очередную доработку, а уточнение архитектуры и ответственности между слоями системы.

Короткие ветки от внешних агентов или сервисных участников тоже имеют смысл, но уже другой. По ним удобно видеть, какие проблемы проект выносил в отдельные срезы: расхождения в \texttt{compare}-workflow, ограничение задержек, правки загрузчиков моделей, уборка устаревшего кода. Для основной повествовательной линии они слишком узкие, зато в аналитике хорошо показывают, какие именно болевые точки становились настолько заметными, что заслуживали отдельной ветки.

\section{К чему пришёл проект и куда он движется}

По истории коммитов видно, что проект прошёл как минимум четыре уровня зрелости.

\begin{enumerate}
\item \textbf{Прототипный уровень.} На раннем этапе команда быстро собирала систему из нескольких источников, приводила к единому виду backend и интерфейс, восстанавливала потоковую обработку и закрывала наиболее болезненные поломки.
\item \textbf{Функциональный уровень.} С появлением \texttt{v3.0}, \texttt{Chainlit}, аппаратной адаптации и новых workflow проект научился решать более широкий класс задач и делать это в зависимости от контекста.
\item \textbf{Платформенный уровень.} В марте система стала выглядеть как настоящий инженерный контур: со state management, registry-подходом, политиками запуска, офлайн-развёртыванием, операторским интерфейсом и наблюдаемостью.
\item \textbf{Консолидированный продуктовый уровень.} В апреле команда начала убирать двусмысленность и закреплять основной интерфейс, а это уже работа не над отдельной функцией, а над устойчивой продуктовой траекторией.
\end{enumerate}

Если говорить простым языком, то главный прогресс проекта заключается не только в количестве написанного кода. Намного важнее то, что история коммитов показывает постепенный переход от режима «сделать, чтобы заработало» к режиму «сделать так, чтобы это можно было сопровождать, развивать и объяснять». В этом смысле текущая линия развития ведёт к более ясной архитектуре, где \texttt{Open WebUI} становится основным пользовательским слоем, а backend, модели, инструменты и миграционный контур получают более чёткие границы ответственности.

Одновременно история показывает и точки внимания. На раннем этапе в репозиторий попадали лишние артефакты окружения, что говорит о ещё неустоявшемся процессе. В отдельные периоды заметна высокая плотность срочных исправлений, а значит, часть архитектурных решений сначала принималась под давлением практических проблем, а уже потом оформлялась в документации и более устойчивых контрактах. Для практики это нормальная картина: именно так часто и выглядит реальная инженерная работа, когда зрелость приходит не в виде одного большого релиза, а через серию корректировок, стабилизаций и переосмыслений.

\section{Заключение}

Анализ истории разработки показывает, что \texttt{llm-tools-platform} за рассматриваемый период прошёл путь от быстрого смешанного прототипа к платформе с выраженным backend- и runtime-контуром. Ключевыми результатами стали переход к версии \texttt{v3.0}, развитие \texttt{RAG} и документных workflow, усиление \texttt{UMS} и управления моделями, формирование операторского слоя и последующая консолидация вокруг \texttt{Open WebUI-first}.

Для меня как для автора отчёта этот разбор важен не только как описание уже сделанного. Он показывает логику развития проекта: сначала создавался работающий каркас, затем добавлялись специализированные возможности, после этого усиливалась операционная устойчивость, а на текущем этапе идёт наведение архитектурной ясности. Такой вывод полезен и для ретроспективы, и для планирования следующего этапа работ, потому что по истории уже видно, какие решения оказались стратегическими, а какие были временными переходными мерами.

\section*{Источники и материалы}
\addcontentsline{toc}{section}{Источники и материалы}

\begin{thebibliography}{9}
\bibitem{repo-readme}
README проекта \texttt{llm-tools-platform}. Локальный файл репозитория: \texttt{README.md}.

\bibitem{repo-tasks}
Операционный backlog и исторические планы проекта. Локальные файлы репозитория: \texttt{TASKS.md}, \texttt{docs/plans/}, \texttt{docs/plans/migration/}.

\bibitem{gost}
НИУ ВШЭ. ГОСТ 7.32-2017 для отчётов о научно-исследовательской работе. URL: \url{https://research.hse.ru/gost}.

\bibitem{nstu}
НГТУ DiSpace. Требования к отчёту и дневнику по практике. URL: \url{https://dispace.edu.nstu.ru/didesk/course/show/11839/4}.

\bibitem{latexgost}
\texttt{latex-g7-32}. Шаблон \LaTeX{} для отчётных документов по ГОСТ. URL: \url{https://github.com/latex-g7-32/latex-g7-32}.

\bibitem{overleaf}
Overleaf. Internship Report Template. URL: \url{https://ru.overleaf.com/latex/templates/internship-report-template/wdhvmvsynxmc}.
\end{thebibliography}

\end{document}
